%% Theory of Computing
%% File: v015a019.tex
%% Publication date: December 21, 2019
%% Authors: James B. Wilson
%
%
%
%

%

%
%

%
%
%
%
%
%
%
%
%
%
%
%
%
%
%
%
%
%
%
%
%
%
%
%

%

\documentclass{toc}

%
%
%
\tocdetails{%
  title = {The Threshold for Subgroup Profiles to Agree is Logarithmic},
%
%
%
%
%
%
  number_of_pages = {25}, 
  number_of_bibitems = {34},
  number_of_figures = {0},
  conference_version = {NONE}, %
    %
    %
%
%
%
  author = {James B. Wilson},
    %
    %
    %
  authorlist = {James B. Wilson},
    %
    %
    %
    %
    %
    %
    %
%
%
  %
  %
  %
  %
%
  acmclassification = {F.2.2},
    %
    %
  amsclassification = {68Q17, 20D15, 20F69, 68Q25},
    %
    %
  keywords = {group isomorphism, profiles, p-groups},
    %
    %
    %
    %
    %
}   %

    %
    %
    %
    %
    %
    %
    %
    %
    %
    %
    %
    %
    %
    %
    %
    %
    %
    %
    %
    %
    %
    %
    %
    %
    %
    %
    %
    %
    %

%
%
\tocdetails{%
volume = 15,
number = 19,
year = 2019,
%
%
%
%
%
%
%
received = {December 20, 2016},
revised = {October 19, 2019},
published = {December 21, 2019},
%
%
doi = {10.4086/toc.2019.v015a019},  %
}   %


%
%
\usepackage{enumerate}

\numberwithin{equation}{section}
\theoremstyle{plain}       %
\newtheorem{thm}[equation]{Theorem}
\newtheorem*{thm*}{Theorem}
\newtheorem{lem}[equation]{Lemma}
\newtheorem{prop}[equation]{Proposition}
\newtheorem{prob}[equation]{Problem}
%
\newtheorem{coro}[equation]{Corollary}
%

\theoremstyle{remark}
%
%

\theoremstyle{definition}
\newtheorem{defn}[equation]{Definition}
\newtheorem{ex}[equation]{Example}  %
\numberwithin{figure}{section}
\numberwithin{table}{section}


\newcommand{\mygroups}{\mathfrak{H}}
\renewcommand{\subset}{\subseteq}
\newcommand{\bbF}{\mathbb{F}}  %
\newcommand{\bbZ}{\mathbb{Z}}  %
\DeclareMathOperator{\Aut}{Aut}  %

\begin{document}

\begin{frontmatter}[classification=text]
  
%
%
%
\title{The Threshold for Subgroup Profiles to Agree is Logarithmic}
%

%
%

%
%
%

%
%
%
%
%
%
%
%
%
%
%
%


%
%
%
%

%

\author[wilson]{James B. Wilson\thanks{This research was supported in part by  NSF grant DMS 1620454.}}

%
%
%
%

%
%
%
%
%
%
\begin{abstract}
  For primes $p>2$ and $e>3$ there are at least $p^{e-3}/e$ groups of order
  $p^{2e+2}$ that have equal multisets of isomorphism types of proper subgroups
  and proper quotient groups, isomorphic character tables, and power maps.
  This
  obstructs recent speculation concerning a path towards efficient isomorphism
  tests for general finite groups. These groups have a special-purpose 
  deterministic polynomial-time isomorphism test.  
\end{abstract}

\iffalse %
%
%
%
%
%
%
%
%
%
%
%
%
%
%
%
%
%
%
%
%
%

\tocarxivcategory{!!!}

\fi %

\end{frontmatter}

%
%
%
\section{Introduction}
The problem \textsc{GraphIso} asks if two graphs of order $n$ are isomorphic.  A
recent breakthrough  by Babai has reduced the complexity of \textsc{GraphIso} to
$n^{O((\log n)^c)}$ for some $c\geq 1$ \cite{Babai:graph}.  An obstacle to
improving this bound is the problem \textsc{GroupIso}, which asks if two finite
groups  given by their multiplication tables are isomorphic.  As implied by work
of Hedrl\'in-Pultr \cite{PH} and shown also by Miller \cite{Miller},
$\textsc{GroupIso}\leq_P \textsc{GraphIso}$.  Amongst the many problems that
reduce to \textsc{GraphIso}, \textsc{GroupIso} generates attention.  See~\cite{ELGOB} for a list of many of the key algorithms over the years.  As a
complexity problem, \textsc{GroupIso} is interesting in part because its
worst-case analysis is largely unchanged from what we obtain by a brute-force
algorithm.  Specifically, by comparing minimal generating sets one can decide
isomorphism against a group $G$ of order $n$ in time $n^{d(G)+O(1)}$ where
$d(G)$ denotes the minimum size of a generating set of a group $G$.  By
Lagrange's theorem $d(G)\leq \log_2 n$ which gives a bound on the complexity of
\textsc{GroupIso} of $n^{\log_2 n+O(1)}$. Rosenbaum explored variations of such
algorithms that lowered the bound to $n^{0.5\log_2 n+O(1)}$~\cite{Rosenbaum}.

A stronger bound on $d(G)$ clarifies the cases of group isomorphism that are
most concerning.  Using the prime factorization $n=p_1^{e_1}\cdots p_s^{e_s}$,
we define the following parameter:
\begin{align}\label{def:mu}
	\mu(p_1^{e_1}\cdots p_s^{e_s}) & =\max\{e_1,\dots, e_s\}.
\end{align}
Guralnick~\cite{Guralnick} proved that $d(G)\leq \mu(n)+1$.  As
$n\to\infty$, the number of integers $n$ for which $\mu(n)\leq c$
tends to $1/\zeta(c+1)$---the reciprocal of the zeta
%
function (see, \eg,~\cite{Stark}).  %
So for 99\% of group orders $n$, group isomorphism is
solved in time polynomial in $n$; 
%
see~\cite{DW:cubefree, DW:most}    %
for an exploration of the complexity of group isomorphism based on the
prime factorization of $n$.

The above suggests that to locate worst-case complexity of group isomorphism we
should consider groups of order $n=p^{\ell}$, $p$ prime, which are known as
\emph{$p$-groups}.  Indeed, families of $p$-groups exist which when given to
present-day group isomorphism tests, \eg, those in \cite{ELGOB}, require
$n^{O(\log_p n)}$ operations.  In light of these difficulties, and the
reductions by Pultr--Hedrl{\'\i}n and Miller, the complexity of $p$-group
isomorphism appears to be a real obstacle to the improvement of graph
isomorphism. 

Recent speculation by Gowers \cite{Gowers:blog} 
%
(see, \eg, Babai \cite[p.   %
693]{Babai:graph}) posed the possibility of using a portion of the subgroups of
a finite group to determine isomorphism types of finite groups.  Isomorphism
tests have been successfully using such ideas as heuristics for some time; 
%
see, \eg, fingerprinting in \cite{ELGOB}.   %
Yet, proving efficiency based on these
heuristics has been obstructed by 90-year-old examples by Rottlaender and
others, that show that subgroup lattices do not characterize groups up to
isomorphism \cite{Rottlaender}.

In \cite{Gowers:blog} Gowers introduced a threshold criterion based on the
following.  We say that a group is \emph{$d$-generated} if $d(G)\le d$. The
\emph{$d$-subgroup profile} of a group $G$ is the multiset of isomorphism types
of $d$-generated subgroups of $G$.  Define $\kappa(G)$ as the minimum $d$ such
that the $d$-profile of $G$ determines the isomorphism type of $G$, and finally
define
\begin{align*}
	\kappa(n) & = \max\{\kappa (G) \mid n=|G|\}.
\end{align*}
Evidently $\kappa(G)\leq d(G)$ and so $\kappa(n)\leq\mu(n)+1$. Comparing
$\kappa(n)$-profiles of groups of order $n$ decides isomorphism in time
$n^{2\kappa(n)+O(1)}$.   Gowers asked for examples where $\kappa(n)\notin O(1)$.
Glauberman and Grabowski \cite{GG} gave instances to show $\kappa(p^{\ell})\in
\Omega(\sqrt{\ell})$. We prove
\begin{thm}\label{thm:main-0}
For a prime $p>2$,\ $\kappa(p^{\ell})\in \Theta(\ell)$. In particular,
isomorphism testing of $p$-groups by comparing subgroup profiles requires
$n^{\Theta(\log_p n)}$ time in the worst case.
\end{thm}

To prove \expref{Theorem}{thm:main-0} we begin with primes $p>2$ and $e>3$ then set
$q=p^e$.  Let $\bbF_q$ be a finite field of order $q$ and define the
\emph{Heisenberg group over $\bbF_q$} as
\begin{align}
	H & = H(\bbF_q) = \left\{
		\begin{bmatrix} 
			1 & \alpha & \gamma \\
			  & 1 & \beta\\
			  & & 1
		\end{bmatrix}
		\;\middle|\;
		\alpha,\beta,\gamma\in \bbF_q\right\}.
\end{align}
Our interest is in the quotients of $H$, which we collect according to their
orders.
\begin{align}
	\mygroups_{p,e}^{g}  = \{G \mid \exists N\lhd H(\bbF_{p^e}), 
		|H(\bbF_{p^e}):N|=p^{2e+g}\textnormal{ and }G\cong H(\bbF_{p^e})/N\}.
\end{align}
We shall see the groups in $\mygroups_{p,e}^{g}$ with $-2e \leq g\leq 0$ are all
elementary abelian and so they are isomorphic to $\mathbb{Z}_p^{2e+g}$.  When
$g=1$ the groups are nonabelian but each is what P.\ Hall called
\emph{extraspecial of exponent $p$} and they are pairwise isomorphic
(\expref{Theorem}{thm:Brahana}). However for $1<g\leq e$ we see increasing
diversity.  Our focus is on $g=2$.


\begin{thm}\label{thm:main}
  $\mygroups_{p,e}^{2}$ has at least $p^{e-3}/e$ isomorphism classes;
  yet, the groups in $\mygroups_{p,e}^{2}$ have equal multisets of
  isomorphism types of proper subgroups, and also of proper quotient
  groups.
\end{thm}

The groups in $\mygroups_{p,e}^{2}$ are an extreme illustration of the
difficulty to capture group isomorphism by a list of invariants.  In
\expref{Section}{sec:invariants} we list numerous other similarities including
character tables, conjugacy classes, and at times automorphism groups.  
%
%
%

Despite these observations, the following class of groups has an efficient
solution to the membership and isomorphism problems.
\begin{align}
	\mygroups_p & = \bigcup_e\bigcup_{1\leq g\leq e} \mygroups_{e,p}^{g}\,.
\end{align}
For computation a group $G$ can be given by numerous methods, most commonly, by
permutations, matrices, or polycyclic presentations;
%
see \cite{Seress}.    %
The complexity of operations in these models differs, so we adopt the
\emph{black-box group} model---as introduced by Babai--Szemer\'edi, see
\cite[Chapter 2]{Seress}. We further assume an oracle for group orders.  This
means we can access compactly encoded generators for the group $G$, we can
compute all group operations, and we can compute $|G|$, but we model these tasks
as requiring only constant time.    In prior work with Lewis, the author proved
the following.

\begin{thm}[\cite{LW}]\label{thm:algo}
Fix a prime $p$.  For groups given as a black-box with an order oracle, there are deterministic polynomial-time algorithms that
\begin{enumerate}[(a)]
\item given a group $G$, decides if $G\in\mygroups_p$; and 
\item given groups $G,\bar{G}\in \mygroups_p$, decides if $G\cong \bar{G}$.
\end{enumerate}
\end{thm}

All permutation groups and their quotients have deterministic polynomial-time
algorithms for the necessary black-box group oracles and to find a group's
order,
%
see~\cite[Chapter 3]{Seress}.    %
Since the rows of the multiplication table
of a group $G$ are a faithful permutation representation of $G$, 
we consider the
multiplication table model as the special case of \emph{regular} permutation
groups. In the context of matrix groups $G$ we assume $p$ is bounded.  In that
case work of Luks implies that we can decide in polynomial time if $G$ is a
$p$-group and if so to compute $|G|$ \cite{Luks:solv-mat}.  So in both these
models we can replace ``black-box polynomial time'' with simply ``polynomial
time.'' The computational complexity of groups given by polycyclic presentations
is nuanced, but in general polynomial-time algorithms for the necessary oracles
are not known, see \cite[Section~4]{collection},~\cite{deep-thought}. Even so
the reports of practical performance with polycyclic groups (\eg,
\cite[Section~7.2]{Seress}) suggest they are equally useful input models.

The algorithms of \cite{LW} have a complexity of $O(p+e^{2\omega}\log^2 p)$,
where $2\leq \omega<3$ is the exponent of matrix multiplication, see
\cite[Chapter~12]{vzG}. In the special case of groups in $\mygroups_{p,e}^{2}$,
Brooksbank, Maglione, and the author give an $O(p^3+e^{\omega}\log^2 p)$-time
algorithm \cite{BMW}.  Note that it takes $\Omega(e^2\log^2 p)$ bits to input
the groups in $\mygroups_{p,e}^{2}$ by any of the standard methods, including
matrices, presentations, or permutations.  These algorithms  have been
implemented in the computer algebra system \textsf{Magma} \cite{Magma} and
decide isomorphism for groups as large as $n=5^{256}$ in about an hour on a
general-purpose personal computer \cite[Figure~1.1]{BMW}. The algorithms in
\cite{BMW} further give a complete invariant for the groups in
$\mygroups_{p,e}^{2}$, which is a homogeneous polynomial of degree $e$ in
$\bbF_p[x,y]$.  This implies that the number of isomorphism classes in
$\mygroups_{p,e}^{2}$ is at most $p^{e}$; 
%
see~\cite[Proposition~9.2]{BMW}.    %
\medskip

Some of the steps in the proof of \expref{Theorem}{thm:main} 
can be extracted as
special cases of results in \cite{LW, BW:autotopism, BMW}. However, there is an
increased need to provide an introduction into the methods in those works.  We
have made this exposition largely self-contained and we rely as much as possible
on proofs based in linear algebra as found in \cite{JacII}.



\subsection*{Preliminaries}
General claims about groups can be found in \cite{Robinson}.  We assume all
groups in this note are finite.  The \emph{Frattini subgroup} $\Phi(G)$ is the
intersection of the maximal subgroups of $G$.  The \emph{exponent} of a group
$G$ is the least positive integer $m$ such that for every $g\in G$,\  $g^m=1$.
The \emph{commutator subgroup} $G'$ is the smallest normal subgroup whose
quotient is abelian, equivalently the subgroup generated by the commutators
$[x,y]=x^{-1}x^y=x^{-1}y^{-1}xy$. $Z(G)$ denotes the \emph{center} of $G$.  We
write $G^p=\langle x^p\mid x\in G\rangle$. As in \cite{BMW}, we borrow
terminology from Lie theory and say the \emph{genus} of a group $G$ of
nilpotence class $2$ and exponent $p$ is
\begin{align}\label{def:genus}
	\textnormal{$p$-genus}(G) & =\log_p |G| - d(G).
\end{align}
We will critically involve the following fact found in \cite[p. 140]{Robinson}.

\begin{thm}[Burnside Basis Theorem]
For a  $p$-group $G$, $\Phi(G)=G' G^p$ and $G/\Phi(G)\cong \mathbb{Z}_p^{d(G)}$.
\end{thm}


We will make frequent use of the fact that $\mathbb{Z}_p^d$ is both an abelian
group and a natural $\bbF_p$-vector space.  In this work all rings, algebras,
and modules are unital (have a unit element), and subrings use the same unit as
the superring.  For emphasis we write $A\leq B$ for subsets which also are
substructures such as subgroups and subalgebras.  Finally we write
$\mathrm{GL}_d(\bbF_q)$ for the group of invertible linear automorphisms on
$\bbF_q^d$, and $\mathrm{SL}_d(\bbF_q)$ for the subgroup of transformations with
determinant 1.


\section{A formula for subgroup profiles}
We prove a formula that, under some hypotheses, calculates the subgroup profiles
in $p$-groups. This allows us to construct groups that produce the same profile
without need to directly compare the groups.

\begin{thm}\label{thm:count}
	Let $G$ be a $p$-group in which $\Aut(G)$ acts transitively on maximal
	subgroups and such that for a maximal subgroup $M$ of $G$,\ $d(G)=1+d(M)$.
	Then for every $J<G$, the size of $\mathcal{J}(J)=\{ K< G \mid K\cong J\}$
	is
	\begin{align*}
		\sum_{f=1}^{1+d(M)} \frac{p^{1+d(M)}-1}{p^f-1} 
		\left|\left\{ K\leq M \;\middle|\;
			K\cong J\textnormal{ and }
			|M:K\Phi(M)|=p^{f-1}\right\}\right|.
	\end{align*}
	Hence the profile map $J\mapsto |\mathcal{J}(J)|$ depends only on the
	isomorphism type of a maximal subgroup of $G$.
\end{thm}

While the hypothesis of \expref{Theorem}{thm:count} is restrictive, in
\expref{Proposition}{prop:max-conj} we demonstrate that the automorphism groups of
groups $G\in \mygroups_{p,e}^{g}$ each contain $\mathrm{SL}_2(\bbF_{p^e})$
acting transitively on the maximal subgroups of $G$, just as
$\mathrm{SL}_2(\bbF_{p^e})$ acts transitively on $\bbF_p$-hyperplanes of
$\bbF_p^{2e}\cong G/\Phi(G)$.  Having targeted this particular action may help
expose the reason we chose Heisenberg groups $H(\bbF_{p^e})$ at the start.
%
%
%
Similar constructions can be fashioned for example from the symplectic group
$\mathrm{Sp}_{2m}(\mathbb{F}_{p^{2e}})$ and other linear subgroups.
%

\begin{lem}\label{lem:Frattini}
Let $G$ be a $p$-group $G$ with a maximal subgroup $M$ having $d(G)=1+d(M)$.
Then $\Phi(G)=\Phi(M)$.
\end{lem}
\begin{proof}
Using the Burnside Basis Theorem on $G$ and on $M$ we calculate 
\begin{align*}
	1 & = \frac{|G:\Phi(G)|\cdot |\Phi(G)|}{|G:M|\cdot |M:\Phi(M)|\cdot |\Phi(M)|}
		=\frac{p^{d(G)} |\Phi(G)|}{ p^{1+d(M)} |\Phi(M)|}
		=\frac{|\Phi(G)|}{|\Phi(M)|}.
\end{align*}
As $\Phi(M)=M' M^p\leq G' G^p=\Phi(G)$, we find that $\Phi(M)=\Phi(G)$.
\end{proof}


\begin{proof}[Proof of \expref{Theorem}{thm:count}]
Fix $J<G$.  We use an $\Aut(G)$-invariant partition
\begin{align*}
	\mathcal{J}(J) & = \bigcup_{f=1}^{d(G)} \mathcal{J}(J,f),
&	\mathcal{J}(J,f) & = 
		\left\{K \in \mathcal{J}(J) \;\middle|\; |G:K\Phi(G)|=p^f\right\}.
\end{align*}
Let $\mathcal{M}$ be the set of maximal subgroups of $G$ and fix $M\in \mathcal{M}$ with $d(G)=1+d(M)$.

%
Fix $f$ and define a bipartite graph between the two sets $\mathcal{J}(J,f)$ and
$\mathcal{M}$, such that $(K,X)\in \mathcal{J}(J,f)\times \mathcal{M}$ is an
edge if, and only if, $K\leq X$.  The action of $\Aut(G)$ on this graph permutes
the vertices of $\mathcal{M}$ transitively.  In particular, the degree of every
vertex $X\in \mathcal{M}$  is the same as the degree of $M$.  Apply
\expref{Lemma}{lem:Frattini} to conclude that $\Phi(G)=\Phi(M)$.  Thus, for every
$K\leq M$, $K\Phi(G)=K\Phi(M)$ and so
\begin{align*}
	{\rm deg}~M & = \left|\left\{ 
			K\leq M \;\middle|\; K\cong J, |G:K\Phi(G)|=p^f 
	\right\}\right|\\
		& = \left|\left\{ K\leq M \;\middle|\; 
		K\cong J,\, |M:K\Phi(M)|=p^{f-1} \right\}\right|.
\end{align*}

Next we compute the degree of $K\in \mathcal{J}(J,f)$, \ie, the
size of the set
\begin{align*}
	\{ X\in \mathcal{M} \mid K\leq X\}
		& = 	\{ X\in \mathcal{M} \mid K\Phi(G)\leq X\}.
\end{align*}
Since $G/\Phi(G)\cong \mathbb{Z}_p^{d(G)}$ and $|G:K\Phi(G)|=p^f$
it follows that
\begin{align*}
	 (G/\Phi(G))/(K\Phi(G)/\Phi(G))\cong 
		\mathbb{Z}_p^{f}.
\end{align*}
In particular the number of maximal subgroups of $G$ containing $K$ equals the number of hyperplanes in an $f$-dimensional $\bbF_p$-vector space.  

At this point we count the number of edges in our graph in two ways.
\begin{align*}
	\frac{p^{d(G)}-1}{p-1}{\rm deg}~M & = 
		\sum_{X\in \mathcal{M}} {\rm deg}~X
		 =	\sum_{K\in \mathcal{J}(J,f)} {\rm deg}~K
		 = |\mathcal{J}(J,f)|\frac{p^{f}-1}{p-1}.
\end{align*}
Thus $|\mathcal{J}(J,f)|= {\rm deg}~M\cdot(p^{1+d(M)}-1)/(p^{f}-1)$.  The theorem follows.
\end{proof}

%
\section{Making \texorpdfstring{$p$}{p}-groups with matrices}\label{sec:making}
We are interested in quotients of groups of $(3\times 3)$-matrices, but it will
be easier to discuss properties of a larger class of groups.  For that we use a
general constructions of $p$-groups that has roots in studies of Brahana and
Baer \cite{Brahana, Baer}.
%
Throughout this section $\bbF=\bbF_p$ is the field of coefficients.

Fix a set $\{T_1,\dots,T_g\}$ of $(r\times c)$-matrices.  Here and throughout
empty blocks in matrices are presumed to be $0$ and $\top$ denotes the
transpose.  Define
\begin{align}\label{def:M}
	M(u,v,w) & :=
	\left[\begin{array}{c|ccc|c}
		1 & & u & & w \\
		\hline
		  & & & & \\	 
		  & & I_r & & 
		   \begin{array}{ccc} T_1v^\top & \cdots & T_g v^\top\end{array} \\
		  & & & &  \\
		 \hline
		  & & & & \\
		  & & & & I_g\\
		  & & & & 
		\end{array}\right]\textnormal{ for } u\in \bbF^r, v\in \bbF^c \textnormal{ and } w\in \bbF^g.
\end{align}
We abbreviate $u*v=(uT_1v^{\top},\ldots,uT_g v^{\top})$ and witness 
multiplication of these matrices behaves as follows.
\begin{align*}
	M(u,v,w)M(u',v',w') & = M(u+u',v+v',
	w+w'+u*v')\,,\\
	[M(u,v,w),M(u',v',w')] & = M(0,0,u*v'-u'*v)\,,\\
	M(u,v,w)^i & = M\left(iu,iv, iw+\binom{i}{2}u *v\right)\,.
\end{align*}
{}From these we fashion the following group of matrices.
\begin{align}\label{eq:Brahana}
	B(T_1,\dots,T_g) & =
		\{ 	M(u,v,w) \mid 
			u\in \bbF^r,	v\in\bbF^c, 
			w\in\bbF^g
	\}.
\end{align}

\begin{prop}\label{prop:basics}
	Assuming $B=B(T_1,\ldots,T_g)$ the following hold.
\begin{enumerate}[(i)]
	\item $B'\leq \{ M(0,0,w)\mid w\in \bbF^g_p\}$ with equality whenever
	$\langle u*v\mid u\in\bbF^r,v\in\bbF^c \rangle=\bbF^g$.
	\item $Z(B)\geq \{ M(0,0,w)\mid w\in \bbF^g_p\}$ with equality if 
	$\bigcap_i \mathrm{Null}(T_i)=0$, where $\mathrm{Null}(T_i)=\{u\mid uT_i=0\}$.
	\item If $p>2$, then $M(u,v,w)^p=1$; hence, $\Phi(B)=B'$ and $\dim \langle T_1,\ldots,T_g\rangle$ is the genus of $B$.
	\item If $\bigcap_i \mathrm{Null}(T_i^{\top})=0$, then $|B|=p^{r+c+g}$.
\end{enumerate}
\end{prop}

Notice if $r,s,g\approx n/3$ then there are $p^{n^3/27+\Theta(n^2)}$ distinct
tuples $T_*=(T_1,\ldots,T_g)$.  Yet $d(B(T_*))\in O(n)$ so the number of
isomorphism classes is bounded by $p^{O(n^2)}$.  As a result these groups
comprise $p^{n^3/27+\Theta(n^2)}$ pairwise distinct isomorphism types of groups
of order $p^n$.  Sims has shown the total number of groups of order $p^n$ is
$p^{2n^3/27+O(n^{2.5})}$ \cite[Chapters~4 \& 5]{BNV:enum}. So despite its humble
appearance, this family is extremely complex. 

As our fields $\bbF_{p^e}$ are finite, there exists an $\omega\in \bbF_{p^e}$
such that $\bbF_{p^e}=\bbF_p(\omega)$.  Hence, $\{1,\omega,\dots,\omega^{e-1}\}$
is a basis for $\bbF_{p^e}$ as an $\bbF_p$-vector space.  Define
$t(\omega)_{ij}^{(k)}\in\bbF_p$ as the \emph{structure constants} so that
\begin{align*}
	\omega^i\cdot \omega^j &= \sum_{k=0}^{e-1} t(\omega)_{ij}^{(k)} \omega^k.
\end{align*}
Also let $T_k\in\mathbb{M}_e(\bbF_p)$ be such that  $[T_{k+1}]_{(i+1)(j+1)}=t(\omega)_{ij}^{(k)}$. 
\begin{ex}
If $\bbF_{p^e}=\bbF_p(\omega)$ then $H(\bbF_q)  \cong B(T_1, \dots, T_e)$.
\end{ex}

In the next section we shall prove the following theorem.
\begin{thm}\label{thm:genus-2}
	For primes $p$ and $e$, $G\in \mygroups_{p,e}^{2}$ if, and only if, there
	are invertible $(e\times e)$-matrices $(T_1,T_2)$ such that $G\cong
	B(T_1,T_2)$ where $T_1^{-1}T_2$ has an irreducible minimal polynomial of
	degree $e$.
\end{thm}

\subsection{Subgroups by row, column, and matrix elimination}
One way to explore the subgroups of the groups $B(T_1,\dots,T_g)$ is to restrict
the range of values of $u$ or $v$ in the formula given in \eqref{eq:Brahana}.
For instance, suppose we restrict the coordinate $u_i=0$.  The result is that
the values in the $i$-th row of each matrix $T_1, \dots, T_g$ can be ignored
within that subgroup.  Hence the subgroup we get is isomorphic to the group we
obtain by first removing the $i$-th row of each matrix in $\{T_1,\dots,T_g\}$
and then using the construction of \eqref{eq:Brahana} to create a group on these
smaller matrices.  Removing one row produces a maximal subgroup, two rows a
subgroup of index $p^2$, and so on.  A similar idea applies to columns.
Reversing the process and inserting rows or columns creates subgroup embeddings.

Restricting values of $w$ in \eqref{eq:Brahana} may result in a subset that is
not closed under multiplication.  A way to avoid that concern is to eliminate
entries $w_i$ only once the corresponding matrix $T_i=0$.  

\begin{ex}
To build an embedding of Brahana groups, \eg, $H(\bbF_3)$ into $H(\bbF_{9})$, we
may add matrices to our list of matrices, or  add to the rows and columns of the
matrices we posses.  In this example we let
\[
  \bbF_9=\left\{ \begin{bmatrix} a_0
      & a_1\\ -a_1 & a_0\end{bmatrix}\;\middle|\; a_0,a_1\in \bbF_3\right\}.
\]
We display the concatenations as partitioned matrices.
\begin{align*}
	H(\bbF_3)  \cong B([1]) & \hookrightarrow B([1],[0]) \\
		& \hookrightarrow B([1 | 0],[0 | 1]) \\
		& \hookrightarrow B\left(
			\begin{bmatrix}
			 1 & 0 \\
			 \hline 
			0 & 1 
			\end{bmatrix},
			\begin{bmatrix} 
			0 & 1 \\ 
			\hline 
			-1 & 0 
			\end{bmatrix}\right)
			\cong H(\bbF_{9}).
\end{align*}
\end{ex}

Using this process in reverse (removing rows, columns, or matrices) lets us
explore many of the subgroups of Brahana groups.  We emphasize that this
approach is not guaranteed to explore every subgroup.
\medskip

In light of \expref{Theorem}{thm:count} we build a family of groups each having a
maximal subgroup of a fixed isomorphism type.  As in \cite[p. 70]{JacII}, for a
polynomial $a(t)=a_0t^0+\cdots +a_{e-1}t^{e-1}+t^e\in \bbF_p[t]$, the
\emph{companion matrix} will be 
\begin{align*}
	C(a(t)) & = 	\begin{bmatrix}
		0 & 1 & \\
		 & \ddots & \ddots\\
		  & & 0 & 1\\
		a_0 & \cdots & & a_{e-1}
	\end{bmatrix}.
\end{align*}

\begin{lem}\label{lem:flat}
	For a polynomial $a(t)$ of degree $e$, the group $G=B(I_e,C(a(t)))$ has a
	maximal subgroup $M$ whose isomorphism type depends only on $p$, $e$, and
	$d(G)=1+d(M)$.
\end{lem}

\begin{proof} 
We delete the last row of $I_e$ and $C(a(t))$ to obtain
\begin{align*}
 	M=B\left(
	\overbrace{\begin{bmatrix}
		1 & 0  & \\
		 & \ddots & \ddots\\
		 & & 1 & 0
	\end{bmatrix}}^e,
	\overbrace{\begin{bmatrix}
		0 & 1 & \\
		 & \ddots & \ddots\\
		 & & 0 & 1
	\end{bmatrix}}^e
	\right)
	\hookrightarrow \qquad \qquad\\
	\qquad\qquad
 	B\left(\begin{bmatrix}
		1 & 0  & \\
		 & \ddots & \ddots \\
		 & & 1 & 0\\
		 \hline
		 0 &\cdots & & 1
	\end{bmatrix},
	\begin{bmatrix}
		0 & 1 & \\
		 & \ddots & \ddots\\
		  & & 0 & 1\\
		  \hline
		a_0 & \cdots & & a_{e-1}
	\end{bmatrix}
	\right)=G.
\end{align*}
Evidently $p^{d(M)}=|M:\Phi(M)|=p^{2e-1}=p^{d(G)-1}$.
\end{proof}

\subsection{Quotient groups by linear combinations}\label{sec:sub-elim}
Next we will want some quotient groups of $B(T_1,\dots, T_g)$.  One can see in
\eqref{eq:Brahana} that for each subset $\{i_1,\dots,i_s\}\subseteq
\{1,\dots,g\}$, there is a natural surjection $B(T_1,\dots,T_g)\to
B(T_{i_1},\dots,T_{i_s})$ and that is indeed a group homomorphism.  This is an
analogue to the way we created subgroups in the previous section.  

Along with reducing the number of matrices we can take linear combinations.
Given scalars $(a_1,\dots,a_g)$ there is an epimorphism from $B(T_1,\dots,T_g)$
to $B(a_1T_1+\cdots +a_g T_g)$.  More generally given a $(g'\times g)$-matrix
$A$, there is an epimorphism
\begin{align*}
	B(T_1,\dots,T_g) & \mapsto B\left(\sum_{j=1}^g A_{1j}T_j,\dots, \sum_{j=1}^g A_{g'j} T_j\right).
\end{align*}


\subsection{Notable isomorphisms}
There are also direct ways to create groups isomorphic to $B(T_1,\dots,T_g)$.
For example, for invertible matrices $X\in\mathbb{M}_r(\bbF_p)$ and
$Y\in\mathbb{M}_c(\bbF_p)$,
\begin{align*}
	B(T_1,\dots,T_g)& \cong B(XT_1Y^{\top},\dots,XT_gY^{\top}).
\end{align*}
We may also permute the order of the matrices.  In particular we can always
insist the first matrix have largest rank and that it be expressed in the form
\[
  \left[\begin{matrix} I_a & 0 \\ 0 & 0 \end{matrix}\right]
\]
through Gaussian elimination.  Thus the groups $B(T_1)$ can be
classified up to isomorphism by the rank of $T_1$.

More generally, for each $i,j\in\{1,\dots,g\}$, and $s\in \bbF_p^{\times}$,
\begin{align*}
	B(T_1,\dots,T_g)& \cong B(T_1,\dots, T_i+s T_j,\dots,T_g) \cong B(T_1,\dots,sT_i,\dots,T_g).
\end{align*}
Thus, if  $\{T'_1,\dots,T'_g\}$ is another basis for $\langle
T_1,\dots,T_g\rangle$,  then $B(T_1,\dots,T_g)\cong B(T'_1,\dots,T'_g)$.  There
can be further isomorphisms between these groups, but these will suffice for our
present discussion.  
\medskip

Using these observations we can make even more complex embeddings.

\begin{lem}\label{lem:b(t)-a(t)}
	For every $a(t),b(t)\in \bbF_p[t]$ with $\deg b(t)=:f<e:=\deg a(t)$, there
	is an embedding $B(I_f, C(b(t)))$ into $B(I_e,C(a(t)))$.
\end{lem}
We emphasize that $b(t)$ has no relation to $a(t)$ other than having lower
degree.  So there is no algebraic reason to guess the possible embedding of
$B(I_f,C(b(t)))$ into $B(I_e,C(a(t)))$.  Manipulating matrices demonstrates how
this is done.
%

\smallskip
\noindent\emph{Proof.}
Fix $b_0,\dots, b_{e-2}\in\bbF_p$.
\begin{align}\label{eq:L1Y}
	\begin{bmatrix} 1 & 0 & & \\ & \ddots & \ddots \\ & & 1 & 0 \end{bmatrix}
	\begin{bmatrix} 1 & & \\ & \ddots & \\ & & 1 \\ b_0 & \dots & b_{e-2} & 1\end{bmatrix}
	& = 	\begin{bmatrix} 1 & 0 & & \\ & \ddots & \ddots &\\ & & 1& 0\end{bmatrix},\\
\label{eq:L2Y}
	\begin{bmatrix} 0 & 1 & & \\ & \ddots & \ddots \\ & & 0 & 1 \end{bmatrix}
	\begin{bmatrix} 1 & & \\ & \ddots & \\ & & 1 \\ b_0 & \dots & b_{e-2} & 1\end{bmatrix}
	& = 	\begin{bmatrix} 0 & 1 & & \\ & \ddots & \ddots &\\ b_0 &\dots & b_{e-2} & 1\end{bmatrix}.
\end{align}
Thus, setting $b(t)=b_0t^0 +\cdots +b_{e-2}t^{e-2}+t^{e-1}$, we obtain the
following embedding.  For the isomorphism we are using the identity
$B(T_1Y^\top,T_2Y^\top)\cong B(T_1,T_2)$ following the calculation of
\eqref{eq:L1Y} and \eqref{eq:L2Y}.
\begin{align*}
	B(I_{e-1}, C(b(t))) & \hookrightarrow
	 B\left(\left[\begin{array}{ccc|c} 1 & 0 & & 0 \\ & \ddots & \ddots &\vdots \\ & & 1& 0\end{array}\right],
	\left[\begin{array}{ccc|c} 0 & 1 & & 0 \\ & \ddots & \ddots & \vdots\\ b_0 &\dots & b_{e-2} & 1\end{array}\right]\right)\\
	&\cong B\left(\begin{bmatrix} 1 & 0 & & \\ & \ddots & \ddots \\ & & 1 & 0 \end{bmatrix},
	\begin{bmatrix} 0 & 1 & & \\ & \ddots & \ddots \\ & & 0 & 1 \end{bmatrix}\right)\\
	\pushQED{\qed} 
	& \hookrightarrow B(I_e, C(a(t))).\qedhere
	\popQED
\end{align*}
%

\begin{remark}
In light of \expref{Theorem}{thm:genus-2}, \expref{Lemma}{lem:b(t)-a(t)} shows that the
groups in $\mygroups_{p,e}^{2}$ each have  $p^{\Omega(e)}$ isomorphism types of
proper subgroups.
\end{remark}

\begin{prop}\label{prop:pairs}
	Given $(e\times e)$-matrices $(T_1,\dots,T_g)$ where $T_1$ is invertible,
	there are polynomials $a_1(t)|\cdots |a_m(t)$ of polynomials in $\bbF_p[t]$,
	and matrices $\tilde{T}_3,\dots,\tilde{T}_g$ such that
\begin{align*}
	B(T_1,\ldots,T_g) & = B(I_e, C(a_1(t))\oplus \cdots \oplus C(a_m(t)), 
		\tilde{T}_3,\dots,\tilde{T}_g).
\end{align*}
\end{prop}
\begin{proof}
Use the Frobenius Normal Form \cite[p. 93]{JacII} to find a divisor chain
$a_1(t)|\cdots |a_m(t)$ and an invertible matrix $X$ such that 
\begin{align*}
	X^{-1}(T_1^{-1}T_2)X& = C(a_1(t))\oplus \cdots \oplus C(a_m(t)).
\end{align*}
Hence,
\begin{align*}
	B(T_1,T_2,\dots,T_g) & \cong 
		B(I_e,T_1^{-1}T_2,\dots,T_1^{-1}T_g) \\
		& \cong B(I_e,C(a_1(t))\oplus \cdots \oplus C(a_m(t)),\dots, 
			X^{-1}T_1^{-1}T_g X).
\end{align*}
So for $3\leq i\leq g$, set $\tilde{T}_i=X^{-1} T_1^{-1}T_2 X$.
\end{proof}

%
\section{Isomorphisms between quotients of Heisenberg groups}
We have so far created many groups and demonstrated ways to construct varied
subgroups and quotients.  Our effort now shifts back to Heisenberg groups and in
particular we will tackle the question of isomorphisms and automorphisms within
$\mygroups_{p,e}^{2}$.  Our main results in this section are proofs of
\expref{Theorem}{thm:genus-2} and the following result.

\begin{thm}\label{thm:adj-tensor}
	Every isomorphism between nonabelian quotients of $H$ of genus $g>1$ lifts
	to an automorphism of $H$.
\end{thm}

This is a special case of \cite[Theorem~4.4]{LW}. We give a self-contained and
largely matrix-based proof.


\subsection{The role of commutation}
A first principle in the theory of nilpotent groups is to treat groups like
algebras by invoking commutation $[x,y]=x^{-1}x^y=x^{-1}y^{-1}xy$ as a
skew-commutative multiplication.  This very nearly distributes over the usual
product, in the following way.
\begin{align}\label{eq:comm-id}
	[xy,z] & = [x,z]^y [y,z], & [x,yz] & = [x,z][x,y]^z.
\end{align}
With $q=p^e$ and $H=H(\bbF_q)$, commutation takes the following form.
\begin{align}\label{eq:comm-H}
	\left[
	\begin{bmatrix}
	1 & \alpha & \gamma\\ & 1 & \beta\\ & & 1 
	\end{bmatrix},
	\begin{bmatrix}
	1 & \alpha' & \gamma'\\ & 1 & \beta'\\ & & 1 
	\end{bmatrix}
	\right] 
	& = 
	\begin{bmatrix}
	1 & 0 & \alpha \beta'-\alpha'\beta\\ & 1 & 0\\ & & 1 
	\end{bmatrix}.
\end{align}
This shows the following two groups are abelian.
\begin{align*}
	H' & = [H,H] =\left\{ 
		\begin{bmatrix} 1 & 0 & \gamma\\ & 1 & 0\\ & & 1 \end{bmatrix} 
	\;\middle|\; \gamma \in\bbF_{q}\right\}, & H/H'\cong \{(\alpha,\beta) \mid 
	\alpha,\beta\in \bbF_q^m\}.
\end{align*}
Evidently there are isomorphisms 
\begin{align}\label{def:iota}
	\iota:H/H'\to \langle \bbF_q^2,+\rangle\textnormal{ and }
	\hat{\iota}:H'\to \langle\bbF_q,+\rangle.
\end{align}
None of these isomorphisms is natural in the category of groups.  In particular
neither $H/H'$ nor $H'$ is an obvious $\bbF_q$-vector space as scalar
multiplication is not part of the operations of a group.  

Normal subgroups are characterized as follows.
\begin{lem}\label{lem:quo}
For $h\in H-H'$, $[h,H]=H'$; thus, if $N$ is normal in $H$ then either $H'\leq N$ or $N\leq H'$. 
\end{lem}

\subsection{Quotients of \texorpdfstring{$H$}{H}}
To inspect the quotients of $H$ we  use a method to ``linearize'' a nilpotent
group which is in some sense the reversal of the constructions we gave in
\expref{Section}{sec:making}.  Early versions of this approach were described by
Brahana and Baer \cite{Baer, Brahana}.

Since elements in $H'=[H,H]$ commute with the whole group, the identities
\eqref{eq:comm-id} imply that commutation factors through $H/H'\times H/H'\to
H'$ and thereby affords a biadditive map 
$[,]_+:\langle \bbF_q^2,+\rangle\times \langle \bbF_q^2,+\rangle\to 
		\langle \bbF_q,+\rangle$,
\begin{align}\label{eq:comm-1}
	[(\alpha,\beta),(\alpha',\beta')]_+ & = \alpha \beta'-\alpha'\beta
		=(\alpha,\beta)\begin{bmatrix} 0 & 1\\ -1 & 0\end{bmatrix} 
			(\alpha',\beta')^{\top}.
\end{align}
To distinguish between the various roles of $[,]$ we let $[,]$ denote group
commutation and $[,]_+$ the biadditive mapping that commutation produces.

\begin{remark}\label{rem:iota}
  The expression in \eqref{eq:comm-1} is misleading.  While
  $\alpha\beta'$ and $\alpha'\beta$ are defined in $\bbF_{q}$, our
  codomain is technically the elementary abelian group
  $\langle \bbF_{q},+\rangle$ and the specific identification depends
  on the perhaps arbitrary choice of maps $(\iota,\hat{\iota})$ of
  \eqref{def:iota}.  Within isomorphism testing such small intuitive
  leaps can hide challenging computations and on a theoretical level
  it remains to be justified that the maps
  $(\iota,\hat{\iota})$---which are not $\bbF_{q}$-linear,
  nevertheless preserve the geometric properties of an
  $\bbF_{q}$-bilinear form.  A more precise formulation of
  \eqref{eq:comm-1} capable of addressing these concerns is given
  below in \eqref{def:gettable}.
\end{remark}

Now \expref{Lemma}{lem:quo} shows that for $N<H'$, $(H/N)'=H'/N$.  So the
commutation of the quotient $H/N$ will accordingly afford a new biadditive map
\begin{align*}
	[,]_+^{H/N}:\langle \bbF_q^2,+\rangle\times \langle \bbF_q^2,+\rangle\to 
		\langle \bbF_q,+\rangle
	\overset{\pi}{\to} \mathbb{Z}_p^g 
\end{align*}
where $\pi$ is given as the homomorphism $\langle \bbF_q,+\rangle\cong H'\to
H'/N\cong \mathbb{Z}_p^g$. The genus of $H/N$ is the value $g$.  

Let us look closely at the case of genus $g=1$.  Fix
$\pi:\langle\bbF_q,+\rangle\to \mathbb{Z}_p$.  Choose a basis
$\{\alpha_1,\dots,\alpha_e\}$ for $\langle \bbF_q,+\rangle$ as an
$\bbF_p$-vector space and such that $\pi(\alpha_i)=1$ if $i=1$ and $0$
otherwise.  Define 
\begin{align*}
	T_{ij}=\pi([(\alpha_i,0),(0,\alpha_j)]) = \pi(\alpha_i \alpha_j).
\end{align*}
Observe 
\begin{align}\label{eq:T-sym}
	T_{ij}=\pi(\alpha_i\alpha_j) =\pi(\alpha_j\alpha_i)=T_{ji}.
\end{align}
So $T=T^{\top}$.  Regarded as a map of $\bbF_p$-vector spaces we see
\begin{align*}
	[(\alpha,\beta),(\alpha',\beta')]_+^{H/N} & = 
		\begin{bmatrix} \alpha & \beta\end{bmatrix}
		\begin{bmatrix} 0 & T\\ -T & 0 \end{bmatrix}
		\begin{bmatrix} \alpha'\\ \beta'\end{bmatrix}.
\end{align*}

As we vary $H/N$ amongst groups of arbitrary genus $1\leq g\leq e$ we describe
$[,]_+=[,]_+^{H/N}$ by a list of linearly independent invertible symmetric
matrices $T_1,\dots,T_g\in \mathbb{M}_e(\bbF_p)$ such that
\begin{align}\label{def:gettable}
	[(\alpha,\beta),(\alpha',\beta')]_+ & = 
		\left(
		\begin{bmatrix} \alpha &\beta \end{bmatrix}\!
		\begin{bmatrix} 0 & T_1\\ -T_1 & 0 \end{bmatrix}\!
		\begin{bmatrix}\alpha'\\ \beta'\end{bmatrix},
		\dots,
		\begin{bmatrix} \alpha &\beta \end{bmatrix}\!
		\begin{bmatrix} 0 & T_g\\ -T_g & 0 \end{bmatrix}\!
		\begin{bmatrix}\alpha'\\  \beta'\end{bmatrix}\right).
\end{align}
This demonstrates the following correspondence.

\begin{thm}[Brahana correspondence]\label{thm:Brahana}
	A group $H/N$ whose commutation is described by linearly independent
	invertible symmetric matrices $(T_1,\dots,T_g)$ where $H/N\cong
	B(T_1,\dots,T_g)$. In particular all quotients $H/N$ of genus $1$ are
	isomorphic.
\end{thm}
\begin{proof}
	The $(T_1,\ldots,T_g)$ are described above.  Recall $M(u,v,w)$ from
	\eqref{def:M}. The required isomorphism $B(T_1,\dots,T_g)\to H/N$ is as
	follows.
\begin{align*}
		M(u,v,w)
	\mapsto
	\begin{bmatrix}
		1 & \iota^{-1}(u) & \hat{\iota}^{-1}(w) \\ & 1 & \iota^{-1}(v)\\ & & 1
	\end{bmatrix}
	\mod{N}.		
\end{align*}
If $g=1$ then $H/N\cong B(T_1)\cong B(I_e)$.  
\end{proof}

\begin{coro}\label{coro:quo-profile}
	The groups in $\mygroups_{p,e}^{2}$ have equal multisets of isomorphism
	types of proper quotient groups.
\end{coro}
\begin{proof}
Fix $N\leq H'$ with $|H:N|=p^{2e+2}$.  As $p^{2e+2}=|H:H'|\cdot
|H':N|=p^{2e}|H':N|$ we find that $|H':N|=p^2$, and so $H/N$ has genus $2$.  As
in \expref{Lemma}{lem:quo}, if $N< K\leq H$ and $K/N$ is normal in $H/N$, then
$K<H'$ or $H'\leq K$.  If $H'\leq K$ then $(H/N)/(K/N)\cong H/K\cong
\mathbb{Z}_p^f$ where $p^f=|H:K|$.  This does not depend on the choice of $N$.
The number of choices for $K$ is the number of subgroups in $\mathbb{Z}_p^{2e}$
of index $f$, which again does not depend on $N$.  Otherwise $N<K<H'$ and so
$|H':K|=p$.  Thus $(H/N)/(K/N)\cong H/K$ has genus $1$.  So by
\expref{Theorem}{thm:Brahana} its isomorphism type is  fixed and independent of $N$.
Finally, $H'/N\cong \mathbb{Z}_p^2$ so there are exactly $p+1$ choices of $K$
with $N<K<H'$.  This is independent of $N$.
\end{proof}

\subsection{Distributive products}\label{sec:tensor}
Throughout this section $\bbF=\bbF_p$ and $q=p^e$.  To prove
\expref{Theorem}{thm:adj-tensor} we need a brief detour to discuss distributive
products.  Take a subset $S\subseteq \mathbb{M}_{r}(\bbF)\times
\mathbb{M}_{c}(\bbF)$.  Note that we do not require that $S$ be a subring or
have any algebraic structure, though often $S$ is a generating set for some
interesting algebraic structure; 
%
see \expref{Remark}{rem:no-ring}.     %
Recall $\top$ denotes transpose. Define 
\begin{align}\label{def:otimes-A}
	\bbF^r\otimes_S\bbF^c & := \{X\in \mathbb{M}_{r\times c}(\bbF)\mid 
		\forall (L,R)\in S, L^\top X=XR\}.
\end{align}
We further develop $\bbF^{r}\otimes_S \bbF^c$ into a distributive product.
First $\mathbb{M}_{r\times c}(\bbF)=(\bbF^{r}\otimes_S \bbF^c)\oplus K(S)$ where
\begin{align}
	K(S) & := \langle L^\top X-XR \mid X\in\mathbb{M}_{r\times c}(\bbF), 
		(L,R)\in S\rangle.
\end{align}
The map $\pi_S$ from $\mathbb{M}_{r\times s}(\bbF)$ onto $\bbF^r\otimes_S
\bbF^c$ with kernel $K(S)$ lets us define a distributive \emph{tensor} product
\begin{align*}
\otimes & =\otimes_S:\bbF^r\times \bbF^c\to \bbF^r\otimes_S\bbF^c,
&
	u\otimes v & = \pi_S(u^\top v).
\end{align*}
Notice for $(L,R)\in S$, $\pi_S(L^\top X)=\pi_S(XR)$ and so we find
\begin{align*}
	(uL)\otimes v & = \pi_S(L^\top (u^\top v))=\pi_S((u^\top v)R)=u\otimes (vR).
\end{align*}

\begin{ex}\label{ex:tensor-M2}
	Let $C\in \mathbb{M}_e(\bbF)$ such that $K=\{ \sum_i a_i C^i \mid a_i\in
	\bbF\}\leq \mathbb{M}_e(\bbF_p)$ is a field of order $p^e$,  \eg, take $C$
	to be the companion matrix of an irreducible polynomial of degree $e$.  Then
	consider the following set of operators in $\mathbb{M}_{2e}(\bbF)\times
	\mathbb{M}_{2e}(\bbF)$.
	\begin{align*}
		S & =\left\{\left(\begin{bmatrix} I_e & 0\\ 0 & C \end{bmatrix},
			\begin{bmatrix} C^{\top} & 0\\ 0 & I_e \end{bmatrix}\right),
			\left(\begin{bmatrix} 0 & I_e \\ I_e & 0 \end{bmatrix},
			\begin{bmatrix} 0 & -I_e\\ -I_e & 0 \end{bmatrix}\right)\right\}.
	\end{align*}
	By solving the system of linear equations of \eqref{def:otimes-A} we find
	\begin{align*}
		\bbF^{2e}\otimes_S \bbF^{2e}
	& = \left\{\begin{bmatrix} 0 & \alpha\\ -\alpha^{\top} & 0 \end{bmatrix} 
			\;\middle|\; \alpha\in K\right\}= \mathrm{Span}_K\left\langle 
			\begin{bmatrix} 0 & I_e\\ -I_e & 0\end{bmatrix}\right\rangle.
	\end{align*}
	In particular $\bbF_p^{2e}\cong K^2$ and $\bbF_p^{2e}\otimes\bbF_p^{2e}\cong
	K$, and $\otimes_S:K^2\times K^2\to K$ is the same as the mapping in
	\eqref{eq:comm-1}.
\end{ex}

Turning now to a general distributive product $\circ:\bbF^r\times \bbF^c\to
\bbF^g$, we can ask if $\circ$ is related to $\otimes_S$.  We do so by comparing
codomains, specifically we define
\begin{align}\label{def:factor-through}
	\otimes_S \to \circ & :\equiv
	(\exists \tau:\bbF^r\otimes_S \bbF^c\to \bbF^g)( u\circ v= \tau(u\otimes_S v) ).
\end{align}
We say $\circ$ \emph{factors through} $\otimes_S$; 
%
see~\cite[(2.4)]{BW:autotopism}.    %

\subsection{Adjoint algebras of distributive products}\label{sec:adjoint}
\expref{Section}{sec:tensor} imagined we possessed a set $S$ of operators from which
we craft a product.  However our circumstance is reversed.  We have distributive
products $[,]_+$ arising from  groups $G\in\mygroups_{p,e}^g$, and these
products have no associated operators.  Even so, we have seen hints, \eg,
\eqref{def:gettable}, that when we identify  $G$ with a quotient of $H(\bbF_q)$,
then the operators $\bbF_q$ seem applicable to the analysis of $[,]_+$.  We
cannot expect to recover something arbitrary like $\bbF_q$ from $[,]_+$, rather
we must set our sights on some form of \emph{universal operators} for $[,]_+$,
as those will persist independent of representation.  We turn to a device
introduced in \cite{Wilson:unique-cent} known as the \emph{adjoint algebra} of a
product.

\begin{defn}
	Fix a distributive product $\circ:\bbF^r\times \bbF^c\to \bbF^g$.  An
	\emph{adjoint pair} $(L,R)\in \mathbb{M}_r(\bbF)\times\mathbb{M}_s(\bbF)$
	satisfies the following for each $u\in \bbF^r$ and each $v\in \bbF^c$
	\begin{align*}
		(uL)\circ v=u\circ (vR).
	\end{align*}
	The \emph{adjoint algebra} $\mathrm{Adj}(\circ)$ is the set of all adjoint
	pairs of $\circ$.
\end{defn}

As the name suggests, $\mathrm{Adj}(\circ)$ is an algebra, \ie, an $\bbF$-vector
space with an $\bbF$-bilinear associative unital multiplication.  However, the
multiplication is not coordinatewise but requires a twist as follows
\begin{align}
	(L_1,R_1)(L_2,R_2) &= (L_1L_2,R_2R_1).
\end{align}
The composition in the second component is also referred to as
\emph{op}-multiplication.  This makes $\mathrm{Adj}(\circ)$ a subring of what is
commonly denoted $\mathbb{M}_{2e}(\bbF)\times \mathbb{M}_{2e}(\bbF)^{\rm op}$.
Indeed $\mathrm{Adj}(\circ)$ is a unital subalgebra---all our rings, modules
and algebras to follow are unital.  Given structure constants for $\circ$, such
as those given in \eqref{def:gettable} for the products $[,]_+$, then we can
compute a basis for $\mathrm{Adj}(\circ)$.

As intended, $\mathrm{Adj}(\circ)$ is a universal set of operators in the
following sense.  Recall \eqref{def:factor-through}.

\begin{thm}[Adjoint-tensor Galois correspondence
{\cite[Theorem~2.11]{BW:autotopism}}] 
\label{thm:Galois}
	Fix a distributive product $\circ:\bbF^r\times\bbF^c\to \bbF_p^g$ and
	$S\subseteq \mathbb{M}_r(\bbF)\times \mathbb{M}_s(\bbF)$.  Then
	\begin{align*}
		 S\subseteq {\rm Adj}(\circ) \Longleftrightarrow 
		 \otimes_S \to \circ.
	\end{align*}
	In particular we obtain two closure operations: $S\subset
	\mathrm{Adj}(\otimes_S)$ and $\otimes_{\mathrm{Adj}(\circ)}\to \circ$.
	Furthermore
	\begin{align*}
		\mathrm{Adj}(\circ) & =\mathrm{Adj}(\otimes_{\mathrm{Adj}(\circ)}), &
		\otimes_S = \otimes_{\mathrm{Adj}(\otimes_S)}.
	\end{align*}
\end{thm}

Notice if $A$ is subalgebra of $\mathbb{M}_r(\bbF)\times \mathbb{M}_s(\bbF)^{\rm
op}$ generated as an algebra by a set $S$, then 
\begin{align*}
	\bbF^r\otimes_{\mathrm{Adj}(\otimes_S)}\bbF^c & 
	\subseteq \bbF^r\otimes_A\bbF^c  \subseteq 
	\bbF^r\otimes_S\bbF^c 
	= \bbF^r\otimes_{\mathrm{Adj}(\otimes_S)}\bbF^c.
\end{align*}
That is, $\bbF^r\otimes_S\bbF^c=\bbF^r\otimes_A\bbF^c
=\bbF^r\otimes_{\mathrm{Adj}(\otimes_S)}\bbF^c$.
In this way one sees there is no advantage to assuming at the start that we have
an algebra over which we tensor.

\begin{remark}\label{rem:no-ring}
  Historically Whitney introduced tensor products as requiring rings, and that
 assumption has been carried into the literature, for example
 \cite[Section~IV.5]{Hungerford}. We have demonstrated that the ring assumption
 has no actual role in defining tensor products and it is often a barrier to
 calculations.  Indeed, even if we intend to tensor over a ring $A$, it is
 often
 the case that $A$ is a proper subring of $\mathrm{Adj}(\otimes_A)$.  So to
 insist on a fixed ring $A$ is arbitrary. It would be enough to regard the
 image
 of $A$ as a set and appeal to $\mathrm{Adj}(\otimes_A)$ for any necessary
 ring  theory.  The main point is $\bbF^r\otimes_S\bbF^c =
 \bbF^r\otimes_{\mathrm{Adj}(\otimes_S)}\bbF^c$ and the left-hand side
 is a more flexible construction.

There are two optional assumptions we have made for simplicity but which may
be removed. The first is to use matrices instead of constructing 
$U\otimes_S V$ as a quotient of some infinite abelian group.  
Secondly, it would be
enough to consider functions 
$S\to \mathbb{M}_r(\bbF)\times \mathbb{M}_c(\bbF)$ instead of subsets.
\end{remark}

One helpful observation is the ability to identify extension of scalars without
prior assumptions.

\begin{lem}\label{lem:centroid}
	Let $K\leq \mathbb{M}_e(\bbF)$ be a subalgebra.
	If $\circ:K^r\times K^c \to K^g$ is biadditive and furthermore satisfies
	\begin{align*}
		&(\forall \alpha\in K) & (\alpha u)\circ v = \alpha(u\circ v) 
			= u\circ(\alpha v),
	\end{align*}
	$u\circ K^c=0$ only if $u=0$, and $K^r\circ v=0$ only if $v=0$;
	then for every $(L,R)\in {\rm Adj}(\circ)$,
	\begin{align*}
		(\alpha u)L& =\alpha(uL), & (\alpha v)R &= \alpha(vR).
	\end{align*}
\end{lem}
\begin{proof} The proof is a so-called ``three pile shuffle'' of the operators:
	\begin{align*}
		((\alpha u)L)\circ v = (\alpha u)\circ (uR)
		=\alpha (u\circ (vR)) 
		= \alpha( (uL)\circ v)
		= (\alpha (uL))\circ v.
	\end{align*}
	As this is for each $v$ and $((\alpha u)L-\alpha (uL))\circ K^c=0$, we get
	that $(\alpha u)L=\alpha (uL)$.  Do likewise with $R$.
\end{proof}

\begin{ex}\label{ex:Hei-Adj}
	Let $\bbF_{p^e}=\bbF_p(\omega)$ and $C=C(\omega)$ be the matrix of left
	multiplication matrix of by $\omega$ on $\langle \bbF_{p^e},+\rangle\cong
	\bbF_p^e$. For $H=H(\bbF_{p^e})$, $[,]_+^H$ as defined in \eqref{eq:comm-H}
	with $N=1$, and $\bbF_{p^e}\cong K=\{\sum_i a_i C^i\mid a_i\in\bbF\}\leq
	\mathbb{M}_e(\bbF)$, we calculate that
	\begin{align*}
	{\rm Adj}([,]_+^H) & = \left\{
	\left(\begin{bmatrix} \alpha &\beta \\\gamma & \delta \end{bmatrix},
	\begin{bmatrix} \delta & -\beta\\ -\gamma & \alpha \end{bmatrix}\right) 
		\;\middle|\; \alpha,\beta,\gamma,\delta\in K
	\right\}\cong \mathbb{M}_2(\bbF_q).
	\end{align*}
	Furthermore, $\langle \bbF_{p^e},+\rangle\cong \langle \bbF_{p^e}^2,
	+\rangle\otimes_{{\rm Adj}([,]_+^H)}\langle \bbF^2_{p^e},+\rangle$.
\end{ex}
\begin{proof}
Consider the matrices $(T_1,\ldots,T_e)$ of \eqref{eq:comm-H}.  Notice from
their definition that for each $\alpha\in K$
\begin{align}
	\alpha T_i & = [(\alpha,0),(0,1)]_+= [(0,1),(\alpha,0)]_+^{-1}
		= T_i\alpha.
\end{align}
In particular $CT_i=T_iC$ for each $i$.  Hence, $S$ from
\expref{Example}{ex:tensor-M2} is contained in ${\rm Adj}([,]_+^H)$.  So, by the
Adjoint-Tensor Galois Correspondence (\expref{Theorem}{thm:Galois}),
\[
\otimes_{S}\to
\otimes_{{\rm Adj}([,]_+^H)}.
\]
That means there is a
$\tau:K\cong \bbF_p^{2e}\otimes_S\bbF_p^{2e}\to \langle \bbF_q,
+\rangle=H'$ which is surjective because of
\expref{Proposition}{prop:basics}(i).  By dimensions this is an
($\bbF_p$-linear) isomorphism.  In particular
${\rm Adj}([,]_+^H)={\rm Adj}(\otimes_S)$.  {}From
Example~\ref{ex:tensor-M2}, $\otimes_S:K^2\times K^2\to K$ satisfies
the hypothesis of \expref{Lemma}{lem:centroid} so
${\rm Adj}(\otimes_S)\subset \mathbb{M}_2(K)\times \mathbb{M}_2(K)$.
Furthermore, letting
\[
J=\begin{bmatrix} 0 & 1 \\ -1 & 0 \end{bmatrix},
\]
 and
$(L,R)\in {\rm Adj}(\otimes_S)$, $L^{\top}J=JR$ so
$R=J^{-1}L^{\top}J$.  Therefore,
${\rm Adj}(\otimes_S)=\{(L,J^{-1}L^{\top}J)\mid L\in
\mathbb{M}_2(K)\}$.
\end{proof}

\subsection{Adjoint algebras of Heisenberg group commutation}

Now we refocus on the goal of \expref{Theorem}{thm:adj-tensor}. 
\begin{lem}\label{lem:Jac}
	Fix a unital $\bbF_p$-subalgebra $K\leq \mathbb{M}_e(\bbF_p)$ where $K$ is a
	field of order $p^e$.  If $A\leq \mathbb{M}_{e}(\bbF_p)$ is a unital
	$\bbF_p$-subalgebra containing $K$, and $e$ prime, then $A=K$ or
	$\mathbb{M}_e(\bbF_p)$.
\end{lem}
\begin{proof}
  Let $V=\bbF_p^e$ be the natural (left) module for $\mathbb{M}_e(\bbF_p)$. So
$V$ is also an $A$-module and an $F$-vector space.  Note that $V$ is
$1$-dimensional as a $K$-vector space.  As $K\leq A$, it follows that for every
$0\neq v\in V$, $V=Kv=Av$.  So $V$ is a simple $A$-module; hence, $A$ is
faithfully represented on a simple module $V$, that is $A$ is primitive.  By the
Wedderburn-Artin theorem, $A$ is therefore isomorphic to $\mathbb{M}_f(\Delta)$
for a finite division ring $\Delta$ \cite[p. 421]{Hungerford}. By the
Dickson-Wedderburn theorem $\Delta\cong \bbF_{p^s}$ for some $s$ \cite[p.
462]{Hungerford}. Lastly, $e=\dim_{\bbF_p} V=s\dim_{\bbF_{p^s}} V=fs$.  As $e$
is prime either $f=1$ and $A=K$, or else $f=e$ and $s=1$ which makes
$A=\mathbb{M}_e(\bbF_p)$.  
\end{proof}

\begin{lem}\label{lem:adj}
If $H/N$ has genus $g>1$ then ${\rm Adj}([,]_+^H)={\rm Adj}([,]_+^{H/N})\cong
\mathbb{M}_2(\bbF_q)$.
\end{lem}
\begin{proof}
We start by observing some necessary adjoints. By the adjoint-tensor Galois
correspondence, ${\rm Adj}([,]_+^H)\leq {\rm Adj}([,]_+^{H/N})$.  Now we claim
this is an equality (as subalgebras of $\mathbb{M}_{2e}(\bbF_p)\times
\mathbb{M}_{2e}(\bbF_p)^{\rm op}$).


We know that the commutation in $H/N$ is given by a set $\{T_1,\dots,T_g\}$ of
linearly independent invertible symmetric matrices \eqref{eq:T-sym}.  So the
linear equations to solve to describe ${\rm Adj}([,]_+^{H/N})$ are the
following.  For each $1\leq i\leq g$,
\begin{align*}
	\begin{bmatrix}
		L_{11} & L_{12}\\ 
		L_{21} & L_{22}
	\end{bmatrix}^{\top}
	\begin{bmatrix}
	0 & T_i\\
	-T_i & 0 
	\end{bmatrix}
	& =
	\begin{bmatrix}
	0 & T_i\\
	-T_i & 0 
	\end{bmatrix}
	\begin{bmatrix}
		R_{11} & R_{12}\\ 
		R_{21} & R_{22}
	\end{bmatrix}.
\end{align*}
For $i=1$ we get 
\begin{align*}
	\begin{bmatrix}
		R_{11} & R_{12} \\
		R_{21} & R_{22}
	\end{bmatrix}
	& =
	\begin{bmatrix}
		0 & T_1 \\
		-T_1 & 0
	\end{bmatrix}^{-1}
	\begin{bmatrix}
		L_{11} & L_{12}\\
		L_{21} & L_{22}
	\end{bmatrix}^{\top}
	\begin{bmatrix}
		0 & T_1 \\ -T_1 & 0
	\end{bmatrix}.
\end{align*}
Now $i=2\leq g$ adds the further constraints $L_{ij}(T_1^{-1}T_2)=(T_1^{-1}
T_2)L_{ij}$. (Here we take advantage of the assumption that the $T_i$ are
symmetric but this condition can be relaxed at the expense of a more complex
formula at this stage.)

Now consider the subalgebra $C(T_1^{-1}T_2) = \{ L\in \mathbb{M}_e(\bbF_p) 
\mid LT_1^{-1} T_2=T_1^{-1} T_2L \}$.  Then
\begin{align*}
	{\rm Adj}([,]_+^{H/N}) & \leq \left\{
	\left(\begin{bmatrix} L_{11} & L_{12}\\ L_{21} & L_{22} \end{bmatrix},
		\begin{bmatrix}
			0 & T_1 \\
			-T_1 & 0
		\end{bmatrix}^{-1}
	\begin{bmatrix} L_{11} & L_{12}\\
		 L_{21} & L_{22} \end{bmatrix}^{\top}
		 \begin{bmatrix}
			0 & T_1 \\
			-T_1 & 0
		\end{bmatrix}\right) 
		 \;\middle|\;
	L_{ij}\in C(T_1^{-1}T_2)\right\}\\
	& \cong \mathbb{M}_2(C(T_1^{-1}T_2).
\end{align*}
By restricting to the first components of ${\rm Adj}([,]_+^{H})\leq {\rm
Adj}([,]_+^{H/N})$, and invoking \expref{Example}{ex:Hei-Adj}, we obtain
$\mathbb{M}_2(K)\leq \mathbb{M}_2(C(T_1^{-1}T_2))$ and by first restricting to
the upper right block we have $K\leq C(T_1^{-1}T_2)\leq \mathbb{M}_e(\bbF_p)$.
If $C(T_1^{-1} T_2)=\mathbb{M}_e(\bbF_p)$ then $T_1^{-1} T_2$ commutes with
every matrix and thus $T_1^{-1} T_2$ is a scalar matrix.  However, $T_2$ and
$T_1$ are linearly independent.  So $T_1^{-1}T_2$ cannot be scalar.  As a result
$C(T_1^{-1}T_2)\neq \mathbb{M}_e(\bbF_p)$.  By \expref{Lemma}{lem:Jac},
$C(T_1^{-1}T_2)=K$.  That is, 
\begin{align*}
	{\rm Adj}([,]_+^{H/N}) & \leq \left\{
		\left(\begin{bmatrix} \alpha & \beta \\ \gamma & \delta\end{bmatrix},
		\begin{bmatrix} \delta & -\beta\\ -\gamma & \alpha\end{bmatrix}\right)
			\;\middle |\;
		\alpha,\beta,\gamma,\delta\in\bbF_q\right\}={\rm Adj}([,]_+^H).
\end{align*}
So indeed ${\rm Adj}([,]_+^H)={\rm Adj}([,]_+^{H/N})$.
\end{proof}

\begin{proof}[Proof \expref{Theorem}{thm:genus-2}]
	Fix $G\in \mygroups_{p,e}^2$.  By \expref{Theorem}{thm:Brahana} we know there
	are linearly independent symmetric invertible matrices $(T_1,T_2)$ such that
	$G\cong B(T_1,T_2)\cong B(I_e,T_1^{-1}T_2)$.  Let $a(t)$ be the minimal
	polynomial of this new $T_1^{-1}T_2$.  As $T_1$ is independent of $T_2$, we
	know that $T_1^{-1}T_2$ cannot be a scalar matrix and so $a(t)$ has degree
	at least $2$.  We need to show $\deg a(t)=e$.

Now let $C(T_1^{-1}T_2)=\{L \in\mathbb{M}_e(\bbF_p) \mid
LT_1^{-1}T_2=T_1^{-1}T_2L\}$.  Following the calculation of the adjoint ring
above we know that
\[
  {\rm Adj}([,]_+^{H/N})\cong \mathbb{M}_2(C(T_1^{-1}T_2)).
\]
Thus, if $B(T_1,T_2)$ is a quotient of $H$ then $C(T_1^{-1}T_2)\cong \bbF_q$.
Since $\bbF_p[t]/(a(t))\cong \bbF_p[T_1^{-1}T_2]\leq C(T_1^{-1}T_2)$ it follows
that $\bbF_p[T_1^{-1}T_2]$ is a subfield of $\bbF_q$.  As $a(t)$ has degree
greater than $1$ and $\bbF_q$ has no intermediate fields, it follows that
$\bbF_p[T_1^{-1}T_2]=\bbF_q$.  Thus $a(t)$ is an irreducible polynomial of
degree $e$.

Conversely if $T_2$ is conjugate to $C(a(t))$ then
\[
  {\rm
Adj}\left([,]_+^{H/N}\right)\cong \mathbb{M}_2(\bbF_q)\cong {\rm
Adj}([,]_+^{H}).
\]
By the Adjoint-Tensor Galois correspondence, the commutation
in $B(I_e,T_1^{-1}T_2)$ factors through the tensor product over ${\rm
Adj}([,]_+^H)$.  By \expref{Example}{ex:Hei-Adj} that tensor is also the commutation
of $H$.  So $B(T_1,T_2)$ is a quotient of $H$.
\end{proof}


\begin{remark}
  At this point we can sketch the concepts behind the algorithms in
  \expref{Theorem}{thm:algo}(a), which we recall serve to recognize
  when a group $G$, given as a black-box group, is a quotient of a
  Heisenberg group.  A key step is to construct the ring
  ${\rm Adj}([,]_+^G)$ and recognize if it isomorphic to
  $\mathbb{M}_2(\bbF_{p^e})$ for some $e$.  If so then the commutation
  in $G$ naturally factors through the tensor over
  $\mathbb{M}_2(\bbF_{p^e})$ which we saw in
  \expref{Example}{ex:Hei-Adj} gives us the commutation of
  $H(\bbF_{p^e})$.  See \cite{BMW, LW} for details.
\end{remark}



\subsection{Automorphisms of Heisenberg groups}
Now we need to consider the automorphisms of $H$, assuming $p>2$.
Each automorphism is described by three constituents:
\begin{enumerate}
\item a homomorphism $\tau:\langle \bbF_q^2,+\rangle 
\to \langle \bbF_q,+\rangle$,
\item an invertible matrix $\begin{bmatrix}\alpha & \beta\\ \gamma & \delta
\end{bmatrix}$ over $\bbF_q$, and
\item a field automorphism $\alpha\mapsto \bar{\alpha}$ of $\bbF_q$.
\end{enumerate}
The corresponding automorphism is as follows.
\begin{align}\label{eq:Aut}
	\begin{bmatrix}
		1 & \alpha' & \gamma'\\ & 1 & \beta' \\ & & 1 
	\end{bmatrix}
	\mapsto
	\begin{bmatrix}
		1 & \overline{\alpha'\alpha+\beta'\gamma} & 
		\overline{(\alpha\delta-\beta\gamma)\gamma'+\tau(\alpha',\beta')}\\ 
		& 1 & \overline{\alpha'\beta+\beta'\delta} \\ 
		& & 1 
	\end{bmatrix}.	
\end{align}
\begin{thm}\label{thm:Aut-Hei}
Every automorphism of $H$ has the form of \eqref{eq:Aut}.
\end{thm}

Our proof will follow the one give in \cite[Proposition~2.9 \& Theorem~4.1]{LW}.
The automorphisms arising from (1), with parts (2) and (3) the identity, form
the subgroup of so-called \emph{central automorphisms}.  We include them for
completeness but remark that they do not influence our proofs to follow.

\begin{remark}
Statements in the literature about the automorphism groups of the of Heisenberg
groups based on symplectic geometry can be misunderstood.  Most claims concern
$\bbF_p$, $\mathbb{Z}$, $\mathbb{R}$ and $\mathbb{C}$, and in the latter cases
considers only smooth automorphisms. As we cautioned in \expref{Remark}{rem:iota},
in the case of $\bbF_q$, $[,]_+:\langle \bbF_q^2,+\rangle\times \langle
\bbF_q^2,+\rangle \to \langle \bbF_q,+\rangle$  is $\bbF_q$-bilinear but
$[,]:H/H'\times H/H'\to H'$ is only biadditive.  So only the cases of
$\mathbb{Z}$ and $\bbF_p$ are immediate by standard geometric methods.  One must
recover the missing ring structure of the coefficients before appealing to
geometric reasoning.
\end{remark}

We saw in \expref{Example}{ex:Hei-Adj} that the commutation of Heisenberg groups is
actually a special type of distributive product, a tensor product.  This means
instead of acting on a biadditive map we can act on a ring ${\rm
Adj}([,]_+)\cong \mathbb{M}_2(\bbF_q)$.


\begin{thm}[Skolem-Noether {\cite[p. 237]{JacII}}] 
The ring automorphisms of $\mathbb{M}_2(\bbF_q)$ are $Y\mapsto X^{-1}\bar{Y}X$
where $X$ is an invertible $2\times 2$ matrix and $\alpha \to \bar{\alpha}$ is a
field automorphism of $\bbF_q$ applied to each entry of $X$.
\end{thm}
\begin{proof}
First the automorphism $\phi$ will send $\alpha I_2\mapsto \bar{\alpha}I_2$
which gives us the field automorphism $\sigma$.  Replacing $\phi$ with
$\phi(Y^{\sigma^{-1}})$ we now have an $\bbF_q$-linear automorphism.  That
function maps the minimal right ideal
\[
  \left\{\left[\begin{smallmatrix}\alpha &
        \beta \\ 0 & 0\end{smallmatrix}\right] : \alpha,\beta\in \bbF_q\right\}
\]
to
another minimal right ideal, either
\[
  \left\{\left[\begin{smallmatrix} 0 & 0 \\
      \gamma & \delta \end{smallmatrix}\right]: \gamma,\delta\in
  \bbF_q\right\} \qquad \text{or}
\qquad \left\{\left[\begin{smallmatrix}\gamma & \delta \\ \nu \gamma &
      \nu\delta\end{smallmatrix}\right]:\gamma,\delta\in
  \bbF_q\right\},
\]
for some $\nu\in \bbF_q$.  Each of these is a
$2$-dimensional vector space over $\bbF_q$ so that transformation can
be given by an invertible square matrix $X$.
\end{proof}

\begin{lem}\label{lem:aut-H}
There is an epimorphism $\Aut(H)\to \Aut(\mathbb{M}_2(\bbF_q))$.  The kernel
consists of the central automorphisms of $H$.
\end{lem}
\begin{proof}
Let $\phi:H\to H$ be an automorphism.  Since $\phi([h,k])=[\phi(h),\phi(k)]$,
$\phi$ factors through $\bbF_p^{2e}\cong H/H'\to H/H'\cong\bbF_p^{2e}$. So we
let $X$ be the matrix representing that transformation.  Also we let $\hat{X}$
be the matrix describing the restriction of $\phi$ to $\bbF_p^e\cong H'\to
H'\cong\bbF_p^e$.  Notice $(X,\hat{X})$ satisfy
\begin{align*}
	[(\alpha,\beta)X,(\alpha',\beta')X]_+ 
		& = [(\alpha,\beta),(\alpha',\beta')]_+\hat{X}.
\end{align*}
Now take $(L,R)\in {\rm Adj}([,]_+)$.  It follows that
\begin{align*}
	[(\alpha,\beta)X^{-1}LX,(\alpha',\beta')]_+
	%
		&= [(\alpha,\beta)X^{-1}L,(\alpha',\beta')X^{-1}]_+\hat{X}\\
%
		& =[(\alpha,\beta),(\alpha',\beta')X^{-1}RX]_+.
\end{align*}
In this way $\Aut(H)$ acts on ${\rm Adj}([,]_+)\cong \mathbb{M}_{2}(\bbF_q)$.
We saw that commutation in $\Aut(H)$ is the same as the tensor product with
${\rm Adj}([,]_+)$ (\expref{Example}{ex:Hei-Adj}).  So every automorphism of ${\rm
Adj}([,]_+)$ determines an automorphism of $H$.  If $(X,\hat{X})$ is the
identity pair then $\phi$ is a central automorphism by definition.
\end{proof}


\begin{proof}[Proof of \expref{Theorem}{thm:adj-tensor}]
In the Brahana correspondence (\expref{Theorem}{thm:Brahana}) we saw that every
nonabelian quotient $H/N$ is determined up to isomorphism  by the matrices
$(T_1,\dots,T_g)$ which also define $[,]_+^{H/N}$.   Fix an isomorphism
$\phi:H/N_1\to H/N_2$.  Since $(H/N_i)/(H/N_i)'\cong
H/H'\cong\langle\bbF_q^2,+\rangle$, we see $\phi$ determines a matrix
\[
  X=\begin{bmatrix} A & B \\ C & D\end{bmatrix} \in \mathbb{M}_{2e}(\bbF_p).
\]
Using the isomorphism $\tau:\langle \bbF_q,+\rangle\to \langle
\bbF_q^2,+\rangle\otimes_{\mathbb{M}_2(\bbF_q)} \langle \bbF_q^2,+\rangle$
described in \expref{Example}{ex:Hei-Adj}, we define $\Gamma:H\to H$ as follows.  If
$c\in \langle\bbF_q,+\rangle$ and $c\tau = \sum_i (a_i,b_i)\otimes (x_i,y_i)$,
then define
\begin{align*}
	\begin{bmatrix}
		1 & a & c \\ & 1 & b \\ & & 1 
	\end{bmatrix}
	& \mapsto
	\begin{bmatrix}
		1 & aA+bC & \sum_i (a_i,b_i)X\otimes (x_i,y_i)X \\ & 1 & aB+bD \\ & & 1 
	\end{bmatrix}.
\end{align*}
{}From our proof of \expref{Lemma}{lem:aut-H} we notice $\Gamma$ is an automorphism of
$H$ if, and only if, $X^{-1}{\rm Adj}([,]_+)X={\rm Adj}([,]_+)$.  Since $\phi$
is an isomorphism $H/N_1\to H/N_2$ we know that
\begin{align*}
	X^{-1}{\rm Adj}\left([,]_+^{H/N_1}\right)X
	={\rm Adj}\left([,]^{H/N_2}_+\right).
\end{align*}
By \expref{Lemma}{lem:adj} we know ${\rm Adj}\left([,]_+\right)={\rm
Adj}\left([,]_+^{H/N_i}\right)$.  So $\Gamma\in \Aut(H)$ and $\Gamma(N_1)=N_2$
and induces $\phi$.
\end{proof}


\begin{coro}\label{coro:count}
The set $\mygroups_{p,e}^{2e+2}$ has at least $p^{e-3}/e$ isomorphism types.
\end{coro}
\begin{proof}
The number of subgroups $N<H'$ of index $p^2$ is the number of subspaces of
codimension 2 in an $\bbF_p$-vector space $H'\cong \bbF_p^e$.  That number is
\[
  \frac{(p^e-1)(p^e-p)}{(p^2-1)(p-1)}.
\]
Meanwhile the action by $\Aut(H)$ on
$H'$ has order $e(p^e-1)$; see \eqref{eq:Aut}.  So the number of orbits is at
least $p^{e-3}/e$.  By \expref{Theorem}{thm:adj-tensor} these orbits are in
bijection with isomorphism types so the result follows.
\end{proof}

\begin{remark}\label{rem:general-iso}
  The proof of \expref{Theorem}{thm:adj-tensor} and the analysis in
  \expref{Corollary}{coro:count} mimics how we efficiently test isomorphism of
  quotients of Heisenberg groups as announced in \expref{Theorem}{thm:algo}(b).  We
  transform the question into one of asking if two $\bbF_p$-subspaces of
  $\bbF_{p^e}$ are in the same orbit under
  $\bbF_{p^e}^{\times}\rtimes {\rm Gal}(\bbF_{p^e})$.  
  As discovered by R{\'o}nyai 
%
  (see~\cite[Lemma~4.8]{LW}),  %
 the orbits under $\bbF_{p^e}^{\times}$ can be decided by solving a system of
 linear equations.  The remaining action by ${\rm Gal}(\bbF_{p^e})$ has small
 orbits.
\end{remark}

\section{Proof of \expref{Theorem}{thm:main}}

In \expref{Corollary}{coro:count} we used the fact that the representation of
$\Aut(H)$ on $H'$ induces the group $\bbF_{p^e}^{\times}\rtimes{\rm
Gal}(\bbF_{p^e})$.  Now we use the fact that the kernel of that representation
induces the group ${\rm SL}(2,\bbF_{p^e})$ on $H/H'$.


\begin{prop}\label{prop:max-conj}
	For every $N\leq H'$, $\Aut(H/N)$ acts 
	transitively on the maximal subgroups of $H/N$.
\end{prop}

We need a few observations, compare \cite[p. 135 \& 140]{Robinson}.

\begin{lem}\label{lem:phi-stuff}
	In a group $G$, for a normal subgroup $N$ contained in $\Phi(G)$ (including
	$N=1$), $\Phi(G/N)=\Phi(G)/N$ and the maximal subgroups of $G$ are in
	bijection with those in $G/N$ and with those in $G/\Phi(G)$.
\end{lem}
	
\begin{lem}\label{lem:basis-2}
	For a $p$-group $G$, there is a bijection between the maximal subgroups of
	$G$ and the projective points in $(\bbF_p^{d(G)})^{\top}$---the dual space
	of $\bbF_p^{d(G)}$.  In particular $\mathrm{Aut}(G)$ acts transitively on
	its maximal subgroups if, and only if, the induced action of
	$\mathrm{Aut}(G)$ on $(\bbF_p^{d(G)})^{\top}$ is transitive on projective
	points.
\end{lem}
\begin{proof}
	By the Burnside basis theorem there is an isomorphism $\tau:G/\Phi(G) \to
\bbF_p^{d(G)}$.  Since each maximal subgroup $M$ of $G$ contains $\Phi(G)$,
$M/\Phi(G)$ is maximal in $G/\Phi(G)$ and so $W=\tau(M/\Phi(G))$ is a
hyperplane.  In particular $W$ is identified with
$\{\phi\in(\bbF_p^{d(G)})^{\top} \mid W\phi=0\}$ as a projective point in
$W^{\top}$.
\end{proof}

\begin{lem}\label{lem:first-trans}
	The natural action of $\mathrm{SL}_2(\bbF_{p^e})$ on the dual space
	$(\bbF_{p^e}^{2})^{\top}$ (\ie, the action on column-vectors) is transitive
	on non-zero vectors, and therefore also on projective points.
\end{lem}
\begin{proof}
	We show $(1,0)$ can reach both $(\alpha,\beta)$ and $(0,\alpha)$, provided
	$\alpha\neq 0$, using determinant 1 matrix.
	\begin{align*}
	\begin{bmatrix} \alpha & 0 \\ \beta & \alpha^{-1} \end{bmatrix}
	\begin{bmatrix} 1\\ 0 \end{bmatrix}
	& = \begin{bmatrix}\alpha\\ \beta\end{bmatrix} &
	\begin{bmatrix} 0 & -\alpha^{-1} \\ \alpha & 0 \end{bmatrix}
	\begin{bmatrix} 1\\ 0 \end{bmatrix}
	& = \begin{bmatrix} 0 \\ \alpha\end{bmatrix}. \qedhere
	\end{align*}
\end{proof}

\begin{proof}[Proof of \expref{Proposition}{prop:max-conj}]
	Fix $N\lhd H$ such that $|H:N|=p^{2e+2}$.  Because $N\leq \Phi(H)$, by
	\expref{Lemma}{lem:phi-stuff}, the maximal subgroups of $H$ are in bijection of
	those of $H/N$.  Following \expref{Proposition}{prop:basics} we know
	$H'=\Phi(H)$ and $(H/N)/\Phi(H/N)\cong H/\Phi(H)\cong\bbF_{p}^{2e}$.  
	
	Notice from \expref{Theorem}{thm:Aut-Hei} that there is a subgroup
	\begin{align*}
		C_{\mathrm{Aut}(H)}(H') & = \{ \phi\in \mathrm{Aut}(H)
			\mid (H')\phi=1_{H'}\}.
	\end{align*}
	As $N<H'$, there is a homomorphism
	\begin{align*}
		\rho:C_{\mathrm{Aut}(H)}(H')\to \mathrm{Aut}(H/N).
	\end{align*}
Furthermore, $\rho$ faithfully embeds the automorphisms of $H$ induced by
matrices
\[
  \begin{bmatrix} \alpha & \beta \\ \gamma & \delta\end{bmatrix}
\]
with
determinant $1$, as described in \expref{Theorem}{thm:Aut-Hei}.  In particular there
is a representation of 
	\begin{align*}
		\mathrm{SL}_2(\bbF_{p^e}) \to \mathrm{Aut}(H/N)
	\end{align*}
	with the natural action on $\bbF_{p^e}^2\cong (H/N)/\Phi(H/N)$. By
	\expref{Lemma}{lem:first-trans}, this action is transitive on projective points.
	By \expref{Lemma}{lem:basis-2} it follows that $\Aut(H/N)$ acts transitively on
	the maximal subgroups of $H/N$.
\end{proof}


\begin{proof}[Proof of \expref{Theorem}{thm:main}]
Corollaries~\ref{coro:count} \& \ref{coro:quo-profile} establish that the set
$\mygroups_{p,e}^{2e+2}$  has at  least $p^{e-3}/e$ isomorphism types and that
these groups all have the same multiset of isomorphism types of proper quotient
groups.  Finally, \expref{Proposition}{prop:max-conj} and \expref{Lemma}{lem:flat} allow
us to invoke \expref{Theorem}{thm:count} to conclude the proof.
\end{proof}

\section{Closing remarks}

\subsection{Further comparisons of the groups \texorpdfstring{$H/N$}{H/N}}\label{sec:invariants}


The groups in $\mygroups_{p,e}^{2}$ are an extreme illustration of the
difficulty to capture isomorphism by a list of invariants.  Along with
equivalent profiles we summarize the following further similarities.
\medskip

Each $G\in \mygroups_{p,e}^{2}$ has $G'=Z(G)=\Phi(G)\cong \mathbb{Z}_p^2$.  Also
the groups have isomorphic lattices of normal subgroups.  They have exponent $p$
and the same multiset of conjugacy class sizes (\expref{Lemma}{lem:quo}).


Next we remark that for any fixed $g\leq e$, the groups in $\mygroups_{p,e}^{g}$
have isomorphic character tables \cite[Section~6.2]{LW}. (See
\cite[p. 232]{Robinson}
for definitions and basic implications of character tables.)  In fact a stronger
property holds.  Observe that character table columns are indexed by conjugacy
classes of a group.  Because $(g^{-1}xg)^m=g^{-1}x^m g$ we can speak of the
$m$-th power map as as function on conjugacy classes of $G$.  An often strong
isomorphism invariant of groups is to consider character table isomorphisms that
also preserve all powers on the conjugacy classes.  For instance, the examples
of Rottlaender have isomorphic character tables as well but they are
distinguished when in addition we add the power-map constraint. For a $p$-group
it suffices to consider the $p$-th powers.  Brauer asked if non-isomorphic
groups could have an isomorphism of character tables that also preserves the
powers of conjugacy classes of the group.  Dade \cite{Dade} provided examples
using groups of order $p^7$, class $3$ and exponent $p$.  Our examples are of
class $2$, and they also answer Brauer's question providing an infinite family
with character tables as large as possible amongst groups of exponent $p$.  See
Nenciu \cite{Nenciu} for examples with composite exponent.


The groups in $\mygroups_{p,e}^{g}$ are directly and centrally indecomposable
and with the same algebraic type of indecomposability (an invariant introduced
in \cite[Theorem~4.41]{Wilson:unique-cent} and~\cite[Theorem~8]{Wilson:RemakI},
that links indecomposability to isomorphism types of local commutative rings and
local Jordan algebras). These features are explored in \cite[Section 5]{LW}. 

Finally we argue that automorphisms often agree, though we cannot guarantee that
they always agree.  By \expref{Theorem}{thm:adj-tensor} we have that for $G\in
\mygroups_{p,e}^{2}$,
\begin{align*}
	\mathrm{Aut}(G) &\hookrightarrow 
	\Gamma{\rm L}_2(\bbF_{p^e})\ltimes \mathbb{M}_{2e\times 2}(\bbF_p), 
	& \Gamma{\rm L}_2(\bbF_{p^e})=\mathrm{Gal}(\bbF_{p^e}/\bbF_p)\ltimes 
		\mathrm{GL}_2(\bbF_{p^e}).
\end{align*}
Then in our proof of \expref{Proposition}{prop:max-conj} we observed
\begin{align*}
	C_{\mathrm{Aut}(G)}(G'):={\rm SL}_2(\bbF_{p^e})\ltimes 
	\mathbb{M}_{2e\times 2}(\bbF_p) &\hookrightarrow
		\mathrm{Aut}(G).
\end{align*}
This means the number of choices for $\mathrm{Aut}(G)$ can be bounded by
counting the subgroups of
\begin{align*}
	\Gamma{\rm L}_2(\bbF_{p^e})/\mathrm{SL}_2(\mathrm{F}_{p^e}) & \cong \mathbb{Z}_e \ltimes \mathbb{Z}_{(p^e-1)/2}.
\end{align*}
Furthermore, $C_{\mathrm{Aut}(G)}(G')$ is the kernel of the induced
representation $\mathrm{Aut}(G)\to \mathrm{GL}_2(\bbF_p)$ given by the action on
$G'\cong \mathbb{Z}_p^2$.  In particular,
$\mathrm{Aut}(G)/C_{\mathrm{Aut}(G)}(G')$ embeds into both the subgroups of
$\mathbb{Z}_e \ltimes \mathbb{Z}_{(p^e-1)/2}$ and $\mathrm{GL}_2(\bbF_{p^e})$.
This limits the order of $\mathrm{Aut}(G)/C_{\mathrm{Aut}(G)}(G')$ to divisors
of 
\begin{align*}
	\ell=\gcd\left(p(p-1)(p+1), e(p^e-1)/2\right)\in O(p^3).
\end{align*}
Whenever $e\nmid \ell$, then it follows that the $\ell$-subgroups of
$\mathbb{Z}_e \ltimes \mathbb{Z}_{(p^e-1)/2}$ lie in $\mathbb{Z}_{(p^e-1)/2}$
and are thus unique.  So the total number of choices for $\mathrm{Aut}(G)$ as
subgroups of $\Gamma{\rm L}_2(\bbF_{p^e})$ is at most the number of subgroups of
a cyclic group of order $\ell$, which is $O(\log \ell)\subset O(\log p)$.  By
\expref{Theorem}{thm:main} we know $\mygroups_{p,e}^{2e+2}$ have at least
$p^{e-3}/e$ distinct isomorphism types and only $O(\log p)$ choices for
automorphisms.  So equivalent automorphism groups, as subgroups of
$\mathrm{Aut}(H)$, are common.  

For example when $p=3$ and $e=7$ then $\ell=1$.   So for each $G\in
\mygroups_{3,7}^{2\cdot 7+2}$, 
\[
\mathrm{Aut}(G)=\mathrm{SL}_2(\bbF_{3^7})\ltimes
\mathbb{M}_{2e\times 2}(\bbF_p)
\]
as a subgroup of $\mathrm{Aut}(H(\bbF_{3^7}))$.
Meanwhile there are at least 19 isomorphism types in
$\mygroups_{3,7}^{2\cdot 7+2}$.


\subsection{Complete invariants}
Pivoting to methods that do give information, 
Barnes and Wall~\cite{Barnes-Wall}
show that the subgroup lattice of a nilpotent group of class $2$ and exponent
$p$ determines the isomorphism type of the group (which corrects an errant
remark of the author).  Yet, the groups in $\mygroups_{p,e}^g$ have maximum
sized lattices with $|G|^{\Theta(\log |G|)}$ subgroups,
chains of length $\log_p |G|$ and antichains of length
$|G|^{\Theta(\log_p |G|)}$.  Presently we have no
tools to compare such large lattices.

Since standard isomorphism invariants are proving unhelpful we suggest one
non-standard but effective tool.  The idea of a polynomial invariant was
suggested to the author by R\'onyai~\cite{Ronyai:private}.  The following
version was defined in~\cite{BMW}.

For groups $G\cong B(T_1,T_2)$ in $\mygroups_{p,e}^2$, we define the
\emph{generalized Pfaffian} as 
\begin{align}\label{def:Pfaffian}
	{\rm Pf}(T_1,T_2) & = \det (x_1 T_1+x_2 T_2)
	\in \bbF_p[x_1,x_2].
\end{align}
As $B(T_1,T_2)\cong B(I_2,C(a(t))$, the generalized Pfaffian can in fact be
computed efficiently as the homogenization of $a(t)$.  The lex-least
representative of the orbit of the generalized Pfaffian under the action by
${\rm GL}_2(\bbF_{p})$ is a defining invariant as proved in \cite{BMW}.  
(Note that
that   %
work deals with a far more general problem.)

\begin{proof}[Sketch of Proof of \expref{Theorem}{thm:algo}]
The recognition of the class $\mathfrak{H}_p$ begins by using a suite of now
standard algorithms to determine if a group $G$ is a $p$-group of nilpotence
class $2$ and exponent $p$.  That work depends on many algorithms, notably
by Sims and Luks; see~\cite{Seress}.  
	
The next stage constructs the commutator map $[,]_+^G$ for $G$ and solves a
system of linear equations to compute the adjoint algebra
$\mathrm{Adj}([,]_+^G)$.  Using results of
%
R\'onyai,  %
Ivanyos, Brooksbank-Luks,
and the author we constructively recognize $\mathrm{Adj}([,]_+^G)$ together
with its op-composition structure; 
%
see~\cite{LW, BW:autotopism}.    %
Following
\expref{Theorem}{thm:genus-2} (and more general versions) this leads
to a constructive recognition of the quotients of Heisenberg groups.
The solution of the system of linear equations to compute
$\mathrm{Adj}([,]_+^G)$ dominates the calculation and takes
$O(e^{2\omega}\log^2 p)$ bit operations.
	
Having recognized quotients of Heisenberg groups we turn to their
isomorphism problem.  The first method of \cite{LW} builds on
\expref{Remark}{rem:general-iso} which shows how to set up a relatively small
search problem with in the Galois group of $\bbF_{p^e}$. In each step of the
search we solve a system of linear equations to decide if there is an
isomorphism.  This leads to an $O(e^{2\omega}\log^2 p)$-time isomorphism
test.

For the faster variation we must solve the system of linear equations
without a quadratic blow-up.  For adjoints this is done in work of
Brooksbank and the author.  We avoid the later linear algebra problem that
arises in the search question by using \eqref{def:Pfaffian}.  We compute the
lex-least representative under the action of $\mathrm{GL}_2(\bbF_p)$.  This
leads to a complexity of $O(p^3+e^{\omega}\log p)$.  See \cite{BMW} for
details and references.
\end{proof}

%
%
%
%
%
%
%
%
%
%
%
%
%
%
%
%
%
%
%
%
%
%
%
%
%
%
%
%
%
%
%
%
%
%
%
%
%
%
%
%
%
%
%
%
%
%
%
%
%
%
%
%
%
%
%
%
%
%
%
%
%
%
%
%
%
%
%
%
%
%
%
%
%
%
%
%
%
%
%
%
%

%
%
%
%
%
%
%
%

%



\subsection*{Acknowledgment}
Thanks to Laci Babai \& Gene Luks who asked me about profiles and offered useful
comments on this article, to Bill Kantor for extensive feedback on this article,
and to the referees for catching blunders and improving the exposition.



\bibliographystyle{tocplain}   %

%
%
%
%
\bibliography{v015a019}

%
%
%
%
%
\begin{tocauthors}
\begin{tocinfo}[wilson]
 James B. Wilson\\
 Associate professor\\
 Department of Mathematics\\
 Colorado State University\\
 Fort Collins, Colorado, USA\\
 James\tocdot{}Wilson\tocat{}ColoState\tocdot{}Edu\\   %
 \url{https://wwww.math.colostate.edu/~jwilson/}      %
\end{tocinfo}
\end{tocauthors}

%
%
%
%
%
%
%
%
%
%
%
%
%
%
\begin{tocaboutauthors}
\begin{tocabout}[wilson]  %
  \textsc{James B. Wilson} graduated with his \phd\ in Mathematics from
  the  %
  \href{https://www.uoregon.edu}{University of Oregon} in 2008 where he studied
  with \href{https://en.wikipedia.org/wiki/William_Kantor}{Bill Kantor
  (advisor)}, \href{https://en.wikipedia.org/wiki/Eugene_M._Luks}{Gene Luks},
  and \href{https://pages.uoregon.edu/wright/research.html}{Charley Wright}.
  Before this he spent nearly four years with the
  \href{https://en.wikipedia.org/wiki/Intel_Architecture_Labs}{Intel
  Architecture Labs} and still enjoys questions from industry.  As a sophomore
  he asked an innocent question about group isomorphism to
  \href{https://www.pdx.edu/math/rudy-beyl}{Professor F. Rudy Beyl}, who kindly
  delayed a response and instead gave him a copy of a lovely paper by 
  %
  \href{https://doi.org/10.1007/BF01181188}{Ada Rottlaender}
  %
  to find the answer 
  himself.  This sparked the author's now
  decades-long   %
  obsession with isomorphism in algebra and its complexity.
\end{tocabout}
\end{tocaboutauthors}

\end{document}

%
%
%
%
